\section{Evaluation}\label{sec:eval}

This section reports what can be measured without a training run at scale. It covers
correctness (tests and export parity), data cost, model cost, learner throughput, and, in
\cref{sec:pipeline-eval}, end-to-end cycles of the full pipeline. None of these numbers is
a statement about playing strength.

\paragraph{Setup.}
Unless stated otherwise, measurements used one Windows~11 workstation:
\begin{itemize}
  \item Intel Core i7-1185G7 (4 cores, 8 logical processors), 16\,GB RAM, no GPU;
  \item Python 3.13.7, PyTorch 2.12.1 (CPU build)~\citep{paszke2019pytorch}, ONNX Runtime
        1.26.0~\citep{onnxruntime2021}.
\end{itemize}
A second agent session sometimes ran builds and tests on the same machine at the same
time, so the timings carry noise. We report medians and treat differences under about
20\% as noise.

\subsection{Correctness}\label{sec:eval-correct}

\paragraph{Tests.}
The learner has 27 automated tests: 19 unit tests, 6 trainer and command-line tests, and
2 export tests. They cover every change described in \cref{sec:learner}:
\begin{itemize}
  \item \textbf{Data.} Shard round-trips, including against the Go-written fixture, plus
        window selection keeping exactly the newest rows.
  \item \textbf{Labels.} The ownership frame of converted JSONL, the disjoint game-level
        split, and alignment of board, policy and ownership under all eight symmetries.
  \item \textbf{Loss.} The masked policy loss is zero on fast-search rows while the value
        loss is not.
  \item \textbf{Architectures.} Recovery from a bare state dict for every preset, the
        \code{gpool} network running on $9\times9$, $13\times13$ and $19\times19$, and
        tolerant loading of checkpoints without an ownership head.
  \item \textbf{Training and export.} Warm-start weight equality, \code{--continue} resuming
        at the saved step, distillation from a legacy teacher into a \code{gpool} student,
        and export of a trained \code{gpool} checkpoint with metadata.
\end{itemize}
One test checks a property of the Go writer: in every row of the committed Go-written
fixture shard, the side-to-move score equals the sum of the ownership map minus komi for
Black to move, and plus komi for White. This holds exactly, which confirms that the Go
writer and the Python reader agree on the side-to-move frame. The orchestrator and Go
test counts are reported in \cref{sec:pipeline-eval}.

\paragraph{Export parity.}
Across 16 random positions, the largest absolute difference between PyTorch and ONNX
Runtime outputs was $2.4\times10^{-7}$ for a trained \code{gpool} checkpoint and
$2.9\times10^{-6}$ for a legacy 4$\times$32 checkpoint from the v3 era. Its architecture
was inferred from tensor shapes, and it lacked an ownership head. Both are far below the
$10^{-3}$ failure threshold. The int8 export loaded and ran, but it is not
parity-checked, since quantisation changes outputs by design.

\subsection{Data cost}\label{sec:eval-data}

We packed the repository's committed self-play file (\code{training/data/samples.jsonl}:
1{,}722 positions from heuristic self-play) into a compressed shard with the learner's
packer:
\begin{itemize}
  \item \textbf{Size.} The file shrank from 4{,}459{,}476 bytes to 116{,}935 bytes, a
        factor of $38.1\times$. The factor for Go-written shards is given in
        \cref{sec:selfplay}.
  \item \textbf{Loading.} The shard loaded in 0.04--0.07\,s over four runs. The v3 dataset loaded the JSONL in
        0.29\,s, but it deferred tensor construction to every epoch. The v4 JSONL
        converter takes 0.9\,s, because it also normalises labels and recovers game
        boundaries.
\end{itemize}
At the v3 replay cap of 50{,}000 rows, the same ratio puts the replay at roughly 3.4\,MB
instead of 130\,MB. That difference decides whether the replay fits in a free CI cache and
syncs quickly off a preemptible instance.

\subsection{Model cost}\label{sec:eval-model}

\Cref{tab:eval-latency} gives single-threaded ONNX Runtime latency for each preset at
batch sizes 1 and 16, using freshly initialised weights (latency does not depend on the
weights). Three points stand out:
\begin{itemize}
  \item \code{gpool}-6$\times$64 is not faster than legacy 6$\times$64 despite having
        $2.7\times$ fewer parameters, consistent with their near-equal MAC counts
        (\cref{tab:learner-arch}). Its extra pooling and slicing operations cost a little
        at batch 1.
  \item legacy-4$\times$48 is about $2\times$ faster, matching ADR~0005's earlier
        measurement.
  \item \code{gpool}-10$\times$96 has $3.5\times$ the MACs of 6$\times$64 but only about
        $1.5\times$ the batch-1 latency, which suggests fixed per-call overheads dominate
        small networks on CPU.
\end{itemize}
Self-play throughput on CPU scales roughly with the inverse of these latencies. That is
why the CPU preset keeps the legacy lineage and the \code{gpool}-10$\times$96 preset is
reserved for GPU runs.

\begin{table}[htbp]
\centering
\small
\begin{tabular}{@{}lrrrr@{}}
\toprule
& \multicolumn{2}{c}{batch 1 (ms)} & \multicolumn{2}{c}{batch 16 (ms)} \\
\cmidrule(lr){2-3}\cmidrule(l){4-5}
Preset & median & p10 & median & p10 \\
\midrule
legacy-6$\times$64 & 2.83 & 2.49 & 34.4 & 27.2 \\
legacy-4$\times$48 & 1.30 & 0.87 & 15.2 & 12.1 \\
\code{gpool}-6$\times$64 & 3.34 & 2.70 & 41.1 & 17.8 \\
\code{gpool}-10$\times$96 & 4.18 & 3.61 & 65.1 & 57.6 \\
\code{gpool}-15$\times$128 & 10.07 & 9.03 & 147.2 & 138.1 \\
\bottomrule
\end{tabular}
\caption{ONNX Runtime CPU latency with one intra-op thread and one inter-op thread, over
300 (batch 1) or 100 (batch 16) timed runs after 20 warm-up runs. The large spread for
\code{gpool}-6$\times$64 at batch 16 reflects machine noise from concurrent jobs.}
\label{tab:eval-latency}
\end{table}

\subsection{Learner throughput}\label{sec:eval-throughput}

\code{python -m gofer\_train.bench} trains for 60 steps at batch 64 on the same data in
two ways: the v3 loop (\code{DataLoader} over the JSONL dataset, legacy 6$\times$64) and
the v4 trainer (device-resident data, augmentation, EMA, gradient clipping).
\Cref{tab:eval-throughput} gives the median and range over three runs on a lightly loaded
machine. On this CPU both are compute-bound at the same rate. The v4 trainer's 9\% edge on
legacy 6$\times$64 is inside our noise threshold: earlier runs taken while the other
session was busy measured both trainers at under half these rates. The v4 per-sample
overheads (symmetry gather, EMA update) are therefore not visible. This is a parity
result, not a speed-up. The design removes per-row Python work and
host-to-device copies, which should matter on a GPU, where the network is no longer the
bottleneck, but we have not measured a GPU.

\begin{table}[htbp]
\centering
\small
\begin{tabular}{@{}llr@{}}
\toprule
Trainer & Network & Samples / s (median of 3) \\
\midrule
v3 loop (\code{DataLoader}) & legacy-6$\times$64 & 397 (395--435) \\
v4 trainer (D4 + EMA) & legacy-6$\times$64 & 432 (422--434) \\
v4 trainer (D4 + EMA) & \code{gpool}-6$\times$64 & 385 (384--390) \\
\bottomrule
\end{tabular}
\caption{CPU training throughput at batch 64 over 60 steps (\code{gofer\_train.bench}); median (min--max) of three runs.}
\label{tab:eval-throughput}
\end{table}

\subsection{What the pipeline found: two defects in the engine}
\label{sec:pipeline-found}

The first thing the v4 gate did was refuse to work. Running it on real networks,
every net-versus-net game was won by White in 12--32 moves, at every playout
count we tried. Because the arena alternates colours, the candidate's score was
then mechanically $0.500$ and no gate could ever decide anything. Chasing that
led to two defects in the engine, both older than v4 and both invisible to the
test suite.

\paragraph{The search preferred the opponent's best move.}
Node values are stored from each node's own side to move: the backup negates the
value at every level, and evaluators return the value from the side to move. A
child's mean is therefore the \emph{opponent's} view, and selection must negate
it. It did not. The search maximised the opponent's value, so more search made
the engine weaker.

The project's own reference baseline measures exactly this: 200 games, identical
heuristic evaluators, Black searching 600 playouts against White's 200
(\code{make reproduce-9x9-baseline}). It had been sitting in the repository
recording \textbf{7 Black wins in 200 games}. Re-running the same command after
the fix gives \textbf{133 of 200}, a win rate of $0.665$, or about $+119$ Elo for
tripling the search:

\begin{center}\small
\begin{tabular}{@{}lccc@{}}
\toprule
Black at 600 playouts vs White at 200 & Black wins & mean game length & games under 10 moves \\
\midrule
before (committed baseline) & 7/200 & 12.9 moves & all won by White \\
after  & 133/200 & 20.0 moves & 5, split 3--2 \\
\bottomrule
\end{tabular}
\end{center}

Black wins two thirds of the games while also conceding komi 6.5, which at these
playout counts favours White (\cref{sec:pipeline-found-komi}), so the search
advantage is larger than the win rate alone suggests. The game-length
distribution is the clearer evidence that play is no longer degenerate: the
handful of remaining short games are split between the colours instead of being
uniformly White wins.

\paragraph{The search could not see the end of the game.}
Two consecutive passes end a game in every game loop, but the search did not
treat them as terminal, and scored finished positions with the evaluator instead
of the rules. Passing into a pass therefore looked ordinary, and with komi the
engine walked into lost endings. We added the two-pass rule and exact terminal
scoring, and kept terminal results out of the transposition table, whose key
covers stones but not pass history. Fixing this alone changed no measurable
outcome; the sign error dominated.

\paragraph{A third defect, in the arena itself.}
With both engine fixes in, a 200-game match between two evaluators that are
\emph{literally the same code} still split 126--74 by role. Per-game seeds were
linear in the game index (\code{seed}\,$+2i$, \code{seed}\,$+2i+1$, and
\code{seed}\,$+(i{+}1)k$ for the opening sampler), while the colour swap
alternates on the same parity $i \bmod 2$. Neighbouring seeds draw correlated
openings, so one parity class---which is exactly one role---was systematically
decided. Measured with identical evaluators over 120 games, one parity class won
100\% of its games as White, and \emph{which} class flipped with the seed:

\begin{center}\small
\begin{tabular}{@{}lccc@{}}
\toprule
& \multicolumn{2}{c}{White wins by parity class} & Role split \\
\cmidrule(lr){2-3}
Seed & unswapped & swapped & (baseline--challenger) \\
\midrule
99   & 51/51 (100\%) & 36/49 & 36--64 \\
7    & 26/41 & 40/40 (100\%) & 55--26 \\
2024 & 30/42 & 43/43 (100\%) & 55--30 \\
\midrule
\multicolumn{4}{@{}l@{}}{\emph{after deriving each seed with a mixing function}} \\
99   & 46/53 & 48/55 & 55--53 \\
7    & 55/60 & 46/60 & 51--69 \\
2024 & 46/59 & 42/57 & 55--61 \\
\bottomrule
\end{tabular}
\end{center}

This mattered more than the raw numbers suggest: the gate attributes wins by
role, so a biased role split is a biased promotion decision. Self-play used the
same consecutive seeds and now uses the same mixing.

\paragraph{Fair komi moved with the fix.}
\label{sec:pipeline-found-komi}
Two tests in the repository encoded a "fair komi" of $-1$ for equal heuristic
engines, with a comment explaining that tournament komi was not fair at test
playout counts. That constant had been fitted to the broken search, which ended
games near-empty in about eleven moves and let komi alone decide them. Measured
again after the fix, at 50 playouts over 120--200 games per point, Black scores
$70\%$ at komi 0, $70\%$ at $1.5$, $53\%$ at $2.0$, $38\%$ at $2.5$, $28\%$ at
$3.5$ and $15\%$ at $6.5$. Fair komi for these engines is therefore near $2.0$,
and the tests were updated to use it.

\paragraph{Why nothing caught it.}
The repository's continuous integration runs \code{go test ./... -short}, and
every arena-scale test skips itself in short mode. One test did observe the
symptom---it logged 9 Black wins in 188 games with identical
evaluators---but it asserted only on role balance and explained the colour skew
away in a comment. We added a unit test of the selection rule that runs in short
mode: given two root children with equal priors and visits, selection must
prefer the one whose subtree is worse for the opponent. It is instant and
deterministic, and it fails on the old code by construction, where the search at
200 and 600 playouts could not be told apart from noise in twenty games.

\paragraph{Consequences for this report.}
Every strength number produced before these fixes is void, including the
committed baseline and the gate scores in \cref{sec:pipeline-eval}. No long
training run had been done, so no trained model was lost; the self-play data
generated so far carries policy targets from a search that preferred the
opponent's reply, and is not worth training on.

\subsection{Pipeline end to end}
\label{sec:pipeline-eval}

\paragraph{Tests.}
The Go shard writer has 4 tests:
\begin{itemize}
  \item NPY header layout and alignment;
  \item a round trip of real self-play output through the ZIP container, checking
        \cref{eq:score-identity} on every row;
  \item next-ply policy labelling, including rows after a fast search;
  \item rejection of shards that would mix board sizes.
\end{itemize}
The orchestrator has 23 tests. They replace Go, PyTorch and ONNX Runtime with a fake
executor that writes the same output files the real tools would. They cover:
\begin{itemize}
  \item a three-cycle run with a seeded first champion, an SPRT acceptance that stops before
        the cap, and an SPRT rejection;
  \item a failure injected into training, followed by a resume that neither repeats
        self-play nor counts rows twice;
  \item reuse of cached gate batches after a restart;
  \item skipping training below the row threshold; hold mode; background self-play;
        in-process and sidecar command lines; pruning of old outputs;
  \item the window formula; shard archiving; configuration parsing and validation; the
        Elo and SPRT arithmetic;
  \item publishing: every champion kept in the registry, the alias and latest release
        moving forward, a failed regression match holding back or demoting a champion,
        failed uploads retried, and rollback that renumbers nothing.
\end{itemize}

\paragraph{End-to-end runs.}
We ran the real loop with the smoke configuration on the machine described in
\cref{sec:eval} (an Intel Core i7-1185G7 with 16\,GB of memory). It used the real Go
engine, learner, exporter and arena, with tiny settings: 8 games per cycle, 24/8 playouts
with $p=0.25$, one or two training epochs, and gates of 20 games at 16 playouts. We ran it
once with sidecar inference for three cycles, and once with in-process inference for two.
The inference backend was the only difference; configuration, seeds and code were the
same. \Cref{tab:pipeline-e2e} lists stage times.

\begin{table}[htbp]
\centering
\small
\begin{tabular}{@{}llrrrrrr@{}}
\toprule
Backend & Cycle & Self-play & Train & Export & Gate & Rows & Window \\
\midrule
sidecar    & 1 & 22.8 & 23.0 & 7.7 & 5.7 & 659 & 561 \\
sidecar    & 2 &  0.0 & 17.7 & 8.4 & 8.9 & 663 & 786 \\
sidecar    & 3 & 20.9 & 15.1 & 7.9 & 9.1 & 380 & 901 \\
in-process & 1 & 17.4 & 21.9 & 6.7 & 3.1 & 659 & 561 \\
in-process & 2 &  0.0 & 15.9 & 7.4 & 1.6 & 663 & 786 \\
\bottomrule
\end{tabular}
\caption{Wall-clock seconds per stage in the end-to-end smoke runs, with self-play rows
produced and the replay window used for training. Self-play time is the time the cycle
waited; with overlap, cycle~2's games were already generated during cycle~1. Cycle~1
gates against the heuristic evaluator; later cycles use SPRT against the champion.
Promotion and publishing each took under 0.2\,s.}
\label{tab:pipeline-e2e}
\end{table}

\paragraph{Observations.}
\begin{itemize}
  \item \emph{Overlap.} Cycle~2's self-play wait was at most 0.01\,s in both runs, because
        its games were played in the background during cycle~1's training, export and
        gate. In cycle~3 of the sidecar run, the wait was 20.9\,s. That cycle's games
        started during cycle~2 but took longer than cycle~2's training and gating, so
        overlap hides self-play only up to the length of the other stages.
  \item \emph{Backend.} The same 20-game gate took 1.6\,s in-process and 8.9\,s with
        sidecars. The sidecar figure includes starting two inference servers and waiting
        for them to be healthy, which dominates at this tiny scale.
  \item \emph{Reproducibility.} The same seeds gave identical row counts with both
        backends (659 and 663), including cycle~2, where self-play used the network. This
        is consistent with the per-game seeding of \cref{sec:selfplay-search}, but it is
        not a proof of bit-identical games.
  \item \emph{Gate behaviour.} The smoke networks are deliberately weak: the first seeded
        champion lost all 10 games to the heuristic evaluator. Each later candidate scored
        exactly 10--10 against it. The likelihood ratio stayed near zero ($-0.10$),
        the 20-game cap was reached, and the fallback rule rejected, which is the intended
        outcome for a candidate with no detectable gain.
  \item \emph{Data.} Shards took 21--41\,KB per cycle. The generated shards were loaded by
        the learner without conversion, and the parity-checked export succeeded in every
        cycle.
\end{itemize}
A full smoke run took about 2.5 minutes of wall-clock time for three cycles, including the
engine build.

\paragraph{Gate length and error rates.}
\Cref{tab:orch-gate-oc} in \cref{sec:orch-gate} gives the exact operating characteristics
of the v3 and v4 gates. It is derived analytically, from the implemented decision code,
not measured from games. Its central results are a false-promotion rate of $2.7\%$ for v4
against $8.0\%$ for v3 at zero true gain, and longer v4 gates near zero gain.

\paragraph{Not yet measured.}
We have not run the container build, any GitHub workflow, a GPU run, or a free-runner
training run. No network trained by v4 has been compared with a v3 champion. These gaps
are listed with the other limitations in \cref{sec:limitations}.

